%% Supplementary Information -- Atmospheric Environment
\documentclass[11pt]{article}
\usepackage[T1]{fontenc}
\usepackage{lmodern}
\usepackage[utf8]{inputenc}
\usepackage[a4paper,margin=25mm]{geometry}
\usepackage{graphicx}
\usepackage{amsmath,amssymb}
\usepackage{booktabs}
\usepackage{url}
\usepackage{rotating}
\newcommand*{\ugm}{\ensuremath{\mu\mathrm{g\,m^{-3}}}}
\newcommand*{\pmten}{\ensuremath{\mathrm{PM_{10}}}}
\renewcommand{\thetable}{S\arabic{table}}
\renewcommand{\thesection}{S\arabic{section}}
\emergencystretch=3em
\sloppy

\begin{document}

\begin{center}
{\large\bfseries Supplementary Information}\\[6pt]
{\large A global reanalysis reproduces when, but not where, PM$_{10}$ is high across the dense
monitoring network of an arid emirate}\\[8pt]
Jin Yoo, Jeongyun Kim, Doo-Hee Lee, Ku Kang
\end{center}

\noindent Every value in this document is produced by the scripts named in each section from the
deposited records; the corresponding result file is given in brackets.

\section{The records, the archive and the model series}

The OpenAQ public archive was retrieved for the 19 location identifiers of the network on
13~November 2026 (\texttt{fetch\_openaq\_archive.py}): 6,803 daily files holding 401,693 rows,
of which 98,459 are \pmten{} (29~October 2022 to 20~May 2024); the remainder are nitrogen dioxide
(99,270), sulfur dioxide (99,218), ozone (68,171) and carbon monoxide (36,575). No value is
negative and no row is duplicated.

The 20,818 three-hourly values up to 31~December 2023 analysed in earlier work on this network
are identical to the archived values at the same time stamps, the archive holds no further value
at a three-hourly time stamp in that period, and the station coordinates agree to within
$5\times10^{-5}$ degrees [\texttt{archive\_checks.json}]. Those values were therefore archived
values sampled every third hour; they were not averaged over three hours. The three-hourly record
of the present study extends the same sampling to 20~May 2024 (32,566 values,
\texttt{build\_records\_full.py}).

The CAMS EAC4 surface \pmten{} field was retrieved from the Copernicus Atmosphere Data Store on
14~November 2026 for 25.5--22.5$^{\circ}$N and 51.75--56.25$^{\circ}$E, 1~October 2022 to
31~May 2024, eight analysis times per day (\texttt{fetch\_cams\_eac4\_pm10.py}). The nine
occupied cells have a value at every one of the 4,872 three-hourly times. At the 20,818 time
stamps of the earlier record the retrieved values agree with the model values used there to
within $1.5\times10^{-5}$~\ugm{} (median $1.1\times10^{-6}$) [\texttt{cams\_cells\_check.json}].

Table~\ref{tab:months} lists what each month contributes. From October 2022 to January 2023 the
archive holds 14 of the 24 hours of the day (06--19~UTC); in February, May and July 2023 it holds
one value per day, at 08~UTC; in April and June 2023 it holds 12 hours; from August 2023 it holds
all 24. The daily record therefore uses the value at 08~UTC, the only hour held in every month.

\begin{table}[ht]\centering\small
\caption{Months held by the three-hourly record and by the archive. ``Days with 08~UTC value'' counts
days on which at least one monitor reports at 08~UTC; the network median and the share of
monitor-days above 45~\ugm{} are computed from those values (24-h means ending 08~UTC).}
\label{tab:months}
\resizebox{0.9\textwidth}{!}{\begin{tabular}{lrrrrr}\toprule
& Analysed & \multicolumn{4}{c}{OpenAQ archive} \\ \cmidrule(lr){3-6}
Month & 3-hourly values & Hourly values & Days with 08 UTC value & Network median & Above 45 (\%) \\ \midrule
2022-10 & 175 & 468 & 2 & 81 & 97 \\
2022-11 & 1{,}526 & 4{,}270 & 17 & 66 & 84 \\
2022-12 & 1{,}945 & 5{,}382 & 21 & 55 & 63 \\
2023-01 & 737 & 2{,}221 & 22 & 50 & 60 \\
2023-02 & -- & 40 & 3 & 68 & 78 \\
2023-04 & 76 & 475 & 17 & 60 & 85 \\
2023-05 & -- & 377 & 20 & 100 & 93 \\
2023-06 & 129 & 816 & 23 & 106 & 91 \\
2023-07 & -- & 208 & 11 & 105 & 96 \\
2023-08 & 1{,}013 & 3{,}135 & 7 & 125 & 98 \\
2023-09 & 4{,}033 & 12{,}338 & 30 & 112 & 98 \\
2023-10 & 2{,}478 & 7{,}655 & 21 & 82 & 92 \\
2023-11 & 4{,}250 & 12{,}627 & 27 & 59 & 79 \\
2023-12 & 4{,}456 & 13{,}285 & 30 & 66 & 87 \\
2024-01 & 4{,}532 & 13{,}538 & 31 & 66 & 85 \\
2024-02 & 3{,}666 & 11{,}066 & 25 & 56 & 68 \\
2024-03 & 2{,}141 & 6{,}341 & 12 & 88 & 93 \\
2024-04 & 925 & 2{,}646 & 8 & 79 & 85 \\
2024-05 & 484 & 1{,}571 & 4 & 96 & 93 \\
\bottomrule\end{tabular}
}
\end{table}

\section{Evidence that the archived values are 24-h running means}

\paragraph{Hourly resolution.} Table~\ref{tab:hourly} compares, for each species reported by the
network, the behaviour of the hourly archive series [\texttt{archive\_checks.json}]. For a
trailing 24-h running mean $x$ of hourly concentrations $h$,
\[
x_t - x_{t-1} = \frac{h_t - h_{t-24}}{24},
\]
so that the increment at $t$ and the increment at $t+24$ share the term $h_t$ with opposite
signs and are negatively correlated, whereas the increments of an hourly concentration with a
diurnal cycle are positively correlated at a lag of 24~h. The \pmten{} increments are negatively
correlated at 24~h at all 19 monitors; those of nitrogen dioxide, ozone and carbon monoxide are
positively correlated at every station that reports them, and those of sulfur dioxide at 16 of
19. The \pmten{} series is also far smoother from hour to hour than any of the gases, and 35\%
of its integer values repeat the previous hour.

\begin{table}[ht]\centering\small
\caption{Hourly signature of the archived series, by species. Medians over stations. ACF(1~h):
lag-1~h autocorrelation of the series. ``Unchanged'': share of values equal to the value one
hour earlier. The last four columns refer to the autocorrelation of the hourly increments; the
range at 24~h is over stations.}
\label{tab:hourly}
\resizebox{\textwidth}{!}{\begin{tabular}{lrrrrrrrr}\toprule
& & & & & \multicolumn{4}{c}{Autocorrelation of hourly increments} \\ \cmidrule(lr){6-9}
Species & Stations & Hourly values & ACF(1 h) & Unchanged & 1 h & 12 h & 24 h [range] & Negative at 24 h \\ \midrule
PM$_{10}$ & 19 & 98{,}459 & 0.996 & 35.0\% & $+0.71$ & $+0.06$ & $-0.36$ [$-0.53$, $-0.19$] & 19/19 \\
NO$_2$ & 19 & 99{,}270 & 0.855 & 0.1\% & $+0.13$ & $-0.01$ & $+0.32$ [$+0.12$, $+0.46$] & 0/19 \\
SO$_2$ & 19 & 99{,}218 & 0.871 & 0.7\% & $+0.08$ & $-0.01$ & $+0.06$ [$-0.02$, $+0.60$] & 3/19 \\
O$_3$ & 13 & 68{,}171 & 0.914 & 0.1\% & $+0.26$ & $-0.10$ & $+0.33$ [$+0.17$, $+0.42$] & 0/13 \\
CO & 7 & 36{,}575 & 0.799 & 8.4\% & $+0.04$ & $+0.01$ & $+0.27$ [$+0.10$, $+0.37$] & 0/7 \\
\bottomrule\end{tabular}
}
\end{table}

\paragraph{Alignment of the averaging window.} Table~\ref{tab:fingerprint} gives, for each
monitor, statistics of the analysed three-hourly record [\texttt{averaging\_fingerprint.json}]:
the lag-3~h autocorrelation, the ratio of the variance of three-hourly increments to that of
24-h increments, and the Pearson correlation of the observed series with the reanalysis at the
same instant ($r_\mathrm{inst}$) and with the 24-h mean of the reanalysis placed so that it ends
at the observation time (trailing), is centred on it, or begins at it (leading). The correlation
decreases from trailing to centred to leading at 17 of the 19 monitors, and the trailing window
exceeds the instantaneous comparison at 18. The median diurnal amplitude of the observed series
is 3\% of the mean.

\paragraph{Reporting convention.} The national guideline for calculating the Emirati Air Quality
Index specifies a 24-h rolling average for \pmten{}, requiring at least 18 valid hourly values.
The station metadata held by the archive do not state an averaging interval, and the inference
above was not confirmed with the network operator.

\begin{table}[ht]\centering\small
\caption{Statistics of the analysed three-hourly record by monitor. $N$: values. ACF(3~h):
lag-3~h autocorrelation. Var.\ ratio: variance of three-hourly increments divided by variance of
24-h increments. $r$: Pearson correlation with the reanalysis at the same instant and with its
24-h mean in three alignments. Diurnal: amplitude of the mean diurnal cycle relative to the
mean.}
\label{tab:fingerprint}
\resizebox{0.95\textwidth}{!}{\begin{tabular}{lrrrrrrrr}\toprule
Series & $N$ & ACF(3 h) & Var.\ ratio & $r_\mathrm{inst}$ & $r_\mathrm{trail}$ & $r_\mathrm{centre}$ & $r_\mathrm{lead}$ & Diurnal (\%) \\ \midrule
Bida Zayed & 1{,}073 & 0.981 & 0.049 & 0.31 & 0.56 & 0.45 & 0.38 & 3.5 \\
Khadeeja School & 1{,}127 & 0.993 & 0.047 & 0.38 & 0.30 & 0.32 & 0.31 & 0.9 \\
Liwa & 1{,}083 & 0.980 & 0.060 & 0.29 & 0.51 & 0.46 & 0.37 & 5.4 \\
Hamdan St. & 1{,}081 & 0.977 & 0.026 & 0.11 & 0.11 & 0.09 & 0.06 & 28.1 \\
Mussafah & 1{,}080 & 0.966 & 0.047 & 0.12 & 0.17 & 0.15 & 0.14 & 10.1 \\
Al Ain St. & 1{,}068 & 0.988 & 0.082 & 0.50 & 0.78 & 0.77 & 0.76 & 3.7 \\
Sweihan & 1{,}099 & 0.980 & 0.077 & 0.38 & 0.63 & 0.63 & 0.61 & 1.3 \\
Khalifa City & 1{,}079 & 0.953 & 0.125 & 0.39 & 0.39 & 0.37 & 0.36 & 3.4 \\
Gayathi & 1{,}144 & 0.981 & 0.065 & 0.52 & 0.67 & 0.62 & 0.58 & 3.2 \\
Al Maqtaa & 1{,}083 & 0.962 & 0.127 & 0.33 & 0.41 & 0.38 & 0.37 & 4.3 \\
Ruwais & 1{,}134 & 0.983 & 0.059 & 0.36 & 0.46 & 0.45 & 0.43 & 1.9 \\
Baniyas School & 1{,}108 & 0.988 & 0.052 & 0.50 & 0.75 & 0.67 & 0.59 & 1.0 \\
Habshan S. & 1{,}076 & 0.945 & 0.191 & 0.40 & 0.64 & 0.53 & 0.44 & 1.0 \\
Al Ain Islamic & 1{,}108 & 0.954 & 0.076 & 0.30 & 0.41 & 0.41 & 0.38 & 3.1 \\
Zakher & 1{,}134 & 0.971 & 0.091 & 0.43 & 0.70 & 0.67 & 0.65 & 0.9 \\
Al Qua'a & 1{,}094 & 0.940 & 0.104 & 0.19 & 0.48 & 0.32 & 0.29 & 1.1 \\
Al Tawia & 1{,}096 & 0.983 & 0.087 & 0.50 & 0.79 & 0.78 & 0.77 & 1.1 \\
E11 Road & 1{,}067 & 0.967 & 0.061 & 0.39 & 0.53 & 0.53 & 0.50 & 15.7 \\
Al Mafraq & 1{,}084 & 0.967 & 0.079 & 0.37 & 0.61 & 0.47 & 0.38 & 9.1 \\
\bottomrule\end{tabular}
}
\end{table}

\clearpage
\section{Matched and instantaneous comparison}

With the reanalysis averaged over the 24~h ending at the observation time, the median Pearson
correlation between model and observations over the monitors is 0.50; with the reanalysis taken
at the same instant it is 0.37. The mean bias of a monitor differs between the two treatments by
at most 2.0~\ugm{}, and the ranking of the 19 monitors by bias is almost identical (Spearman
$\rho=0.998$) [\texttt{station\_representativeness.json}].

\section{Stability across sub-periods}

Table~\ref{tab:subperiods} repeats the daily model comparison for four sub-periods
[\texttt{station\_representativeness.json}], and Table~\ref{tab:extended} the statistics that
need no model [\texttt{extended\_observations.json}]. The ranking of the monitors by mean bias in
each sub-period agrees with that of the whole record (Spearman 0.92 to 0.98), the metropolitan
monitors are underestimated by 60 to 69~\ugm{} in each, and the Al~Ain monitors are within
11~\ugm{}. The station means exceed 45~\ugm{} at all 19 monitors in every sub-period. The share of
the between-monitor variance of the station means that lies within model cells is between 42\%
and 51\%, and that of the station medians between 35\% and 44\%.

\begin{table}[ht]\centering\small
\caption{Daily model comparison by sub-period (24-h means ending 08~UTC). Bias in \ugm{} (CAMS
minus observed): range over the 19 monitors and mean over the metropolitan and the Al~Ain
monitors. Rank agreement: Spearman correlation of the monitors' mean bias in the sub-period with
that of the whole record. Regional $\rho$: Spearman correlation of the observed and modelled
network median, days with at least ten monitors.}
\label{tab:subperiods}
\resizebox{\textwidth}{!}{\begin{tabular}{lrrrrrr}\toprule
Sub-period & Days & Bias range & Metropolitan & Al Ain & Rank agreement & Regional $\rho$ \\ \midrule
Oct 2022--Dec 2023 & 251 & $-145$ to $+43$ & $-63$ & $-1$ & 0.98 & 0.70 \\
Sep--Dec 2023 & 108 & $-144$ to $+38$ & $-69$ & $+3$ & 0.92 & 0.71 \\
Outside Sep--Dec 2023 & 223 & $-145$ to $+44$ & $-60$ & $-6$ & 0.97 & 0.65 \\
Jan--May 2024 & 80 & $-144$ to $+39$ & $-64$ & $-11$ & 0.95 & 0.55 \\
\bottomrule\end{tabular}
}
\end{table}

\begin{table}[ht]\centering\small
\caption{Observation-only statistics by sub-period. One value per monitor and day (24-h mean
ending 08~UTC). Pair statistics use pairs with at least 30 common days.}
\label{tab:extended}
\resizebox{\textwidth}{!}{\begin{tabular}{lrrrrr}\toprule
 & To 31 Dec 2023 & Sep--Dec 2023 & Outside Sep--Dec 2023 & Jan--May 2024 & Full archive \\ \midrule
Days & 251 & 108 & 223 & 80 & 331 \\
Monitor-days & 4{,}520 & 2{,}001 & 3{,}988 & 1{,}469 & 5{,}989 \\
Monitor means, range (\ugm{}) & 57--202 & 56--199 & 55--199 & 51--189 & 55--199 \\
Monitor means above 45 & 19/19 & 19/19 & 19/19 & 19/19 & 19/19 \\
Monitor-days above 45 (\%) & 85.2 & 89.1 & 81.8 & 81.4 & 84.3 \\
Variance of monitor means within cells (\%) & 44.8 & 50.7 & 41.8 & 50.1 & 45.7 \\
Variance of monitor medians within cells (\%) & 36.4 & 39.9 & 35.5 & 43.5 & 35.3 \\
Largest within-cell range of means (\ugm{}) & 107 & 104 & 106 & 102 & 105 \\
Day-by-day SD between monitors, median (\ugm{}) & 37.3 & 35.6 & 35.1 & 31.8 & 35.6 \\
Day-by-day SD within cells, median (\ugm{}) & 17.0 & 18.3 & 16.3 & 16.2 & 17.0 \\
Spearman, same-cell pairs (median) & 0.61 & 0.68 & 0.59 & 0.63 & 0.62 \\
Spearman, different-cell pairs (median) & 0.48 & 0.52 & 0.49 & 0.51 & 0.47 \\
Same-cell pair-days on opposite sides of 45 (\%) & 13.5 & 11.2 & 15.8 & 16.7 & 14.2 \\
\bottomrule\end{tabular}
}
\end{table}

\section{Statistical correction: further results}

\paragraph{Effect of the averaging interval.} Table~\ref{tab:three} gives the full set of 24
comparisons under two treatments of the same 32,135 screened three-hourly values: the
instantaneous model value at the observation time, and the mean of the model over the 24~h ending
then [\texttt{audit\_instantaneous.json}, \texttt{audit\_matched.json}]. Under the instantaneous
treatment 16 of the 24 comparisons establish a gain; under matched averaging none does. The two
treatments use identical rows, so the difference is due to the averaging alone. In the earlier
work on this network the record was confined to August--December 2023 for matched averaging and
the same contrast appeared there: a leave-one-station-out gain of $+0.25$ to $+0.38$ under the
instantaneous treatment on the matched rows, and of $+0.10$ to $+0.15$, not established, under
matched averaging [\texttt{v13period/}].

\paragraph{Guard interval.} Table~\ref{tab:guard3} repeats the two stages that use time blocks
with a guard of three days in place of ten [\texttt{audit\_matched\_guard3.json}]. No comparison
differs from zero.

\paragraph{Which predictors carry the gain.} Table~\ref{tab:features} gives the
leave-one-station-out gain of the coordinate-free correction for four sets of predictors
[\texttt{feature\_sets\_loso.json}]. The hour of day adds little; the day of the year carries
the gain in both treatments. The last row reproduces the coordinate-free leave-one-station-out
entry of Table~\ref{tab:three}.

\paragraph{Software.} Python with scikit-learn 1.9.1, NumPy 2.5.3, pandas 3.0.5 and SciPy
1.18.1. The gradient-boosting settings are those of \texttt{audit\_v7.py} in the repository.

\begin{sidewaystable}[p]\centering\footnotesize
\caption{Median $\Delta$, number of monitors improved and Benjamini--Hochberg $q$ for the 24
comparisons under two treatments of the averaging interval, on the same values.}
\label{tab:three}
\resizebox{0.75\textheight}{!}{\begin{tabular}{llrcc}\toprule
Stage & Features & $R$ & Instantaneous model value & Matched 24-h mean \\ \midrule
buffer & with coordinates & 0 & $+0.323$, 17/19, 0.009 & $+0.186$, 13/19, 1.000 \\
buffer & with coordinates & 25 & $+0.185$, 15/19, 0.029 & $+0.001$, 10/19, 1.000 \\
buffer & with coordinates & 50 & $+0.179$, 12/19, 0.411 & $-0.158$, 7/19, 1.000 \\
buffer & with coordinates & 100 & $+0.152$, 11/19, 0.648 & $-0.000$, 9/19, 1.000 \\
buffer & coordinate-free & 0 & $+0.418$, 16/19, 0.015 & $+0.113$, 12/19, 1.000 \\
buffer & coordinate-free & 25 & $+0.260$, 16/19, 0.015 & $-0.000$, 9/19, 1.000 \\
buffer & coordinate-free & 50 & $+0.223$, 16/19, 0.015 & $-0.026$, 9/19, 1.000 \\
buffer & coordinate-free & 100 & $+0.223$, 15/19, 0.029 & $-0.191$, 9/19, 1.000 \\
+ time & with coordinates & 0 & $+0.170$, 17/19, 0.009 & $+0.038$, 11/19, 1.000 \\
+ time & with coordinates & 25 & $+0.108$, 15/19, 0.029 & $-0.111$, 8/19, 1.000 \\
+ time & with coordinates & 50 & $+0.058$, 13/19, 0.200 & $-0.291$, 7/19, 1.000 \\
+ time & with coordinates & 100 & $+0.080$, 11/19, 0.648 & $-0.000$, 9/19, 1.000 \\
+ time & coordinate-free & 0 & $+0.171$, 16/19, 0.015 & $+0.016$, 10/19, 1.000 \\
+ time & coordinate-free & 25 & $+0.118$, 16/19, 0.015 & $+0.016$, 10/19, 1.000 \\
+ time & coordinate-free & 50 & $+0.169$, 15/19, 0.029 & $+0.008$, 10/19, 1.000 \\
+ time & coordinate-free & 100 & $+0.145$, 15/19, 0.029 & $-0.050$, 9/19, 1.000 \\
+ time + cell & with coordinates & 0 & $+0.192$, 14/19, 0.085 & $+0.091$, 11/19, 1.000 \\
+ time + cell & with coordinates & 25 & $+0.099$, 14/19, 0.085 & $-0.111$, 9/19, 1.000 \\
+ time + cell & with coordinates & 50 & $+0.058$, 13/19, 0.200 & $-0.291$, 7/19, 1.000 \\
+ time + cell & with coordinates & 100 & $+0.080$, 11/19, 0.648 & $-0.000$, 9/19, 1.000 \\
+ time + cell & coordinate-free & 0 & $+0.146$, 15/19, 0.029 & $+0.008$, 10/19, 1.000 \\
+ time + cell & coordinate-free & 25 & $+0.118$, 15/19, 0.029 & $+0.008$, 10/19, 1.000 \\
+ time + cell & coordinate-free & 50 & $+0.169$, 15/19, 0.029 & $+0.008$, 10/19, 1.000 \\
+ time + cell & coordinate-free & 100 & $+0.145$, 15/19, 0.029 & $-0.050$, 9/19, 1.000 \\
\bottomrule\end{tabular}
}
\end{sidewaystable}

\begin{table}[p]\centering\footnotesize
\caption{As Table~3 of the main text, with a three-day guard. $q$ is corrected across the 24
comparisons of this run, which include the eight without time blocks.}
\label{tab:guard3}
\resizebox{\textwidth}{!}{\begin{tabular}{rrcrrcrr}\toprule
& & \multicolumn{3}{c}{With coordinates} & \multicolumn{3}{c}{Coordinate-free} \\ \cmidrule(lr){3-5}\cmidrule(lr){6-8}
$R$ (km) & Sources & Median $\Delta$ [95\% CI] & Improved & $q$ & Median $\Delta$ [95\% CI] & Improved & $q$ \\ \midrule
\multicolumn{8}{l}{\emph{Time blocks with three-day guard}} \\
0 & 18.0 & $+0.083$ [$-0.097$, $+0.210$] & 12/19 & 1.000 & $+0.016$ [$-0.260$, $+0.267$] & 10/19 & 1.000 \\
25 & 15.5 & $-0.045$ [$-0.382$, $+0.109$] & 8/19 & 1.000 & $+0.016$ [$-0.276$, $+0.213$] & 10/19 & 1.000 \\
50 & 14.4 & $-0.284$ [$-0.447$, $+0.117$] & 8/19 & 1.000 & $+0.011$ [$-0.282$, $+0.184$] & 10/19 & 1.000 \\
100 & 12.0 & $-0.086$ [$-1.447$, $+0.167$] & 9/19 & 1.000 & $-0.021$ [$-0.370$, $+0.199$] & 9/19 & 1.000 \\
\multicolumn{8}{l}{\emph{+ monitors sharing the model cell withheld}} \\
0 & 16.4 & $+0.086$ [$-0.360$, $+0.282$] & 11/19 & 1.000 & $+0.011$ [$-0.288$, $+0.276$] & 10/19 & 1.000 \\
25 & 15.2 & $+0.057$ [$-0.382$, $+0.168$] & 10/19 & 1.000 & $+0.011$ [$-0.288$, $+0.191$] & 10/19 & 1.000 \\
50 & 14.4 & $-0.284$ [$-0.430$, $+0.117$] & 8/19 & 1.000 & $+0.011$ [$-0.282$, $+0.184$] & 10/19 & 1.000 \\
100 & 12.0 & $-0.086$ [$-1.447$, $+0.167$] & 9/19 & 1.000 & $-0.021$ [$-0.370$, $+0.199$] & 9/19 & 1.000 \\
\bottomrule\end{tabular}
}
\end{table}

\begin{table}[p]\centering\small
\caption{Leave-one-station-out gain (median $\Delta$ and number of monitors improved) of the
correction for four sets of predictors, without coordinates.}
\label{tab:features}
\begin{tabular}{lcc}\toprule
Predictors & Instantaneous CAMS & Matched 24-h mean \\ \midrule
CAMS value & $+0.046$, 12/19 & $-0.199$, 8/19 \\
CAMS value, hour of day & $+0.056$, 12/19 & $-0.215$, 8/19 \\
CAMS value, day of year & $+0.423$, 16/19 & $+0.111$, 12/19 \\
CAMS value, hour of day, day of year & $+0.418$, 16/19 & $+0.113$, 12/19 \\
\bottomrule\end{tabular}

\end{table}

\clearpage
\section{Season, aggregation of the local deficit, and interpolation}

All values in this section use the daily record (24-h means ending 08~UTC) and the days with at
least ten monitors [\texttt{revision\_diagnostics.json}, \texttt{revision\_diagnostics.py}].
The warm season is May to September and the cool season November to February; March, April and
October enter only the analyses over all days.

\paragraph{Season.} The terciles of the regional level are largely seasonal
(Table~\ref{tab:season}, upper part): 77 of the 107 days of the highest tercile are warm-season
days and 94 of the 109 days of the lowest are cool-season days. Terciles formed within each season
(lower part) show the same pattern as those over all days: the metropolitan and the Al~Ain bias
both become more negative as the regional level rises, in each season.

\begin{table}[ht]\centering\small
\caption{Regional terciles by season. Upper part: season of the days in each tercile over all
days. Lower part: terciles formed within each season; regional level is the median over the days
of the tercile; bias (CAMS minus observed, \ugm{}) is the mean over monitor-days of the
metropolitan and the Al~Ain monitors, and the difference is metropolitan minus Al~Ain.}
\label{tab:season}
\begin{tabular}{llrrrrr}
\toprule
Days & Tercile & Days & Regional level & Metropolitan & Al~Ain & Difference \\
\midrule
All days & low & 109 & \multicolumn{4}{l}{cool 94, transition 13, warm 2} \\
All days & mid & 111 & \multicolumn{4}{l}{cool 72, transition 23, warm 16} \\
All days & high & 107 & \multicolumn{4}{l}{cool 10, transition 20, warm 77} \\
\midrule
Cool (Nov--Feb) & low & 61 & 48 & $-$39 & $+$15 & $-$54 \\
 & mid & 57 & 61 & $-$56 & $+$6 & $-$62 \\
 & high & 58 & 74 & $-$63 & $-$12 & $-$51 \\
Warm (May--Sep) & low & 32 & 84 & $-$50 & $+$2 & $-$53 \\
 & mid & 33 & 109 & $-$77 & $-$9 & $-$68 \\
 & high & 30 & 139 & $-$115 & $-$41 & $-$75 \\
\bottomrule
\end{tabular}

\end{table}

\paragraph{Aggregation of the local deficit.} Table~\ref{tab:offset} gives the difference between
the metropolitan and the Al~Ain bias under different aggregations. In the lowest tercile every
aggregation gives between $-50$ and $-62$~\ugm{}. In the highest tercile the value depends on
the aggregation and on Al~Mafraq, the monitor with the largest bias. The daily regression of the
difference on the regional level uses days with at least four metropolitan and two Al~Ain
monitors; intervals are from 10,000 bootstrap resamples of days. No slope differs from zero.

\begin{table}[ht]\centering\small
\caption{Metropolitan minus Al~Ain bias (\ugm{}) by tercile of the regional level, under
different aggregations, and its daily regression on the regional level (slope in \ugm{} per
\ugm{} of regional level).}
\label{tab:offset}
\begin{tabular}{lrrr}
\toprule
Aggregation & Low & Mid & High \\
\midrule
Mean over monitor-days (main text) & $-$59 & $-$57 & $-$63 \\
As above, without Al~Mafraq & $-$52 & $-$49 & $-$41 \\
Mean of per-monitor means & $-$60 & $-$58 & $-$63 \\
Median of per-monitor means & $-$50 & $-$51 & $-$33 \\
One metropolitan monitor left out (range) & $-$62 to $-$52 & $-$61 to $-$49 & $-$69 to $-$41 \\
Terciles within the cool season & $-$54 & $-$62 & $-$51 \\
Terciles within the warm season & $-$53 & $-$68 & $-$75 \\
\midrule
Daily regression on the regional level & Days & Slope & 95\% interval \\
\midrule
All days & 324 & $-$0.00 & $-$0.23 to $+$0.22 \\
Warm season & 95 & $+$0.06 & $-$0.38 to $+$0.39 \\
Cool season & 174 & $+$0.44 & $-$0.58 to $+$1.17 \\
All days, without Al~Mafraq & 323 & $+$0.18 & $-$0.02 to $+$0.39 \\
Warm season, without Al~Mafraq & 95 & $+$0.10 & $-$0.34 to $+$0.40 \\
\bottomrule
\end{tabular}

\end{table}

\paragraph{Interpolation.} Table~\ref{tab:interp} repeats the main statistics with the model
interpolated bilinearly from the four surrounding grid points in place of the nearest one. The
station-level values differ from those of the main text, which use the three-hourly record.

\begin{table}[ht]\centering\small
\caption{Nearest grid point and bilinear interpolation, daily record. Concentrations and biases
in \ugm{}.}
\label{tab:interp}
\begin{tabular}{lrr}
\toprule
Statistic (daily record) & Nearest grid point & Bilinear \\
\midrule
Mean bias, range over monitors & $-$144 to $+$42 & $-$147 to $+$21 \\
Mean bias, mean over monitors & $-$26 & $-$25 \\
Model station mean, range & 42 to 112 & 48 to 91 \\
Spearman, observed vs model station means & $-$0.27 & $-$0.43 \\
Pearson $r$ in time, median over monitors & 0.54 & 0.55 \\
Regional level, Spearman & 0.67 & 0.68 \\
Model SD between monitors, median over days & 23 & 19 \\
Metropolitan minus Al~Ain, low tercile & $-$59 & $-$61 \\
Metropolitan minus Al~Ain, mid tercile & $-$57 & $-$60 \\
Metropolitan minus Al~Ain, high tercile & $-$63 & $-$58 \\
\midrule
\multicolumn{3}{l}{Largest change of a monitor's mean bias: 21; Spearman of the two bias rankings: 0.95} \\
\bottomrule
\end{tabular}

\end{table}

\section{Dust from a sun photometer and from airport weather reports}

AERONET Version~3 daily averages were queried on 14~November 2026 for the 24 sites in the
Emirates, and the record used was retrieved with \texttt{fetch\_aeronet.py}. Over 1~October 2022 to 31~May 2024 only DEWA Research
Centre (24.77$^{\circ}$N, 55.37$^{\circ}$E) holds a long record, 427 days at Level~1.5; Masdar
Institute holds 43 days, Mezaira 13 and the other 21 sites none. The coarse-mode aerosol optical
depth at 500~nm of the spectral deconvolution algorithm, averaged over the day of the observation
and the previous day, is available for 228 days of the daily record
[\texttt{aeronet\_dust\_check.json}]. Its Spearman correlation with the observed regional level is
0.70 and with the modelled one 0.54; the fine-mode fraction correlates with the observed regional
level at $-0.59$. Table~\ref{tab:aeronet} gives the bias by tercile of the coarse-mode optical
depth and the cross-classification with the regional terciles. On the least dusty third of days
the mean bias of the seven metropolitan monitors lies between $-97$ (Al~Mafraq) and $-30$~\ugm{}
(Baniyas School).

\paragraph{Airport weather reports.} The METAR reports of Abu Dhabi International Airport
(OMAA) were retrieved from the archive of the Iowa Environmental Mesonet on 14~November 2026
(\texttt{fetch\_metar\_omaa.py}; 14,632 reports) [\texttt{metar\_dust\_check.json}]. Days need at
least 18 reports in the 24-h window ending 08~UTC; 327 days of the daily record qualify. Dust or
sand is reported 56 times (DU 42, BLDU 14), on 14 days; haze is reported 234 times. On the days
with a dust report the median coarse-mode optical depth at the sun photometer is 0.23, against
0.13 on the others. The share of reports with visibility reduced below 10~km without
precipitation, fog or mist correlates with the observed regional level at 0.35 and with the
coarse-mode optical depth at 0.13, because haze of any origin reduces visibility.
Table~\ref{tab:metar} gives the regional level, the regional bias and the local deficit (medians
over days) by class of day. The difference in local deficit between days with and without dust
reports has $p=0.26$ (Mann--Whitney), that in regional bias $p=0.08$. On the 150 clear days the
mean bias of the metropolitan monitors lies between $-137$ (Al~Mafraq) and $-33$~\ugm{}
(Baniyas School).

\begin{table}[ht]\centering\small
\caption{Regional level, model minus observed regional level and metropolitan minus Al~Ain bias
(medians over days, \ugm{}) by class of day in the weather reports of Abu Dhabi International
Airport.}
\label{tab:metar}
\resizebox{\textwidth}{!}{\begin{tabular}{lrrrrr}
\toprule
Days & $N$ & Regional level & Model minus observed & Metropolitan minus Al~Ain & (without Al~Mafraq) \\
\midrule
Dust or sand reported & 14 & 125 & $-$37 & $-$72 & $-$50 \\
No dust reported & 313 & 69 & $-$21 & $-$58 & $-$48 \\
Clear (no obscuration) & 150 & 64 & $-$15 & $-$60 & $-$48 \\
\bottomrule
\end{tabular}
}
\end{table}

\begin{table}[ht]\centering\small
\caption{Bias by tercile of the AERONET coarse-mode aerosol optical depth at DEWA Research Centre
(228 days). Concentrations and biases in \ugm{}; bias is CAMS minus observed, mean over
monitor-days. Lower part: number of days in each combination of regional and coarse-AOD tercile.}
\label{tab:aeronet}
\begin{tabular}{lrrr}
\toprule
Tercile of coarse-mode AOD & Low & Mid & High \\
\midrule
Days & 76 & 76 & 76 \\
Coarse-mode AOD, median & 0.06 & 0.13 & 0.25 \\
Fine-mode fraction, median & 0.69 & 0.55 & 0.33 \\
Observed regional level, median & 58 & 70 & 102 \\
Modelled regional level, median & 41 & 51 & 70 \\
Model minus observed regional level, median & $-$14 & $-$26 & $-$32 \\
Bias, metropolitan monitors & $-$53 & $-$61 & $-$73 \\
Bias, Al~Ain monitors & $+$7 & $-$3 & $-$20 \\
Metropolitan minus Al~Ain & $-$60 & $-$58 & $-$53 \\
Metropolitan minus Al~Ain, without Al~Mafraq & $-$53 & $-$46 & $-$30 \\
\midrule
Regional tercile (rows) by coarse-AOD tercile & Low & Mid & High \\
\midrule
Regional low & 49 & 23 & 10 \\
Regional mid & 27 & 33 & 15 \\
Regional high & 0 & 20 & 51 \\
\bottomrule
\end{tabular}

\end{table}

\section{Dust, sea salt and wind in the reanalysis}

The total, dust and sea-salt aerosol optical depth at 550~nm and the 10-m wind components of
CAMS EAC4 were retrieved from the Atmosphere Data Store on 14~November 2026 for the same area,
period and analysis times as the \pmten{} field (\texttt{fetch\_cams\_eac4\_aod\_wind.py}). Each
quantity was averaged over the eight analysis times of the 24~h ending 08~UTC and over the three
cells that hold the metropolitan monitors; the wind direction is that of the vector-mean wind
[\texttt{dust\_wind\_analysis.json}]. The difference between the metropolitan and the Al~Ain bias
is computed per day from days with at least four metropolitan and two Al~Ain monitors (324 days).

In the cell of the DEWA Research Centre sun photometer the dust optical depth of the reanalysis
correlates with the AERONET coarse-mode optical depth at 0.85, and its dust share of the optical
depth with the AERONET fine-mode fraction at $-0.76$ (226 days). Over the metropolitan cells dust
accounts for a median 30\% of the modelled optical depth and sea salt for a median 1.5\%
(maximum 9.9\%). Table~\ref{tab:dustwind} stratifies the metropolitan minus Al~Ain difference.

\begin{table}[ht]\centering\small
\caption{Metropolitan minus Al~Ain bias (median over days, \ugm{}; in brackets without
Al~Mafraq) by tercile of the CAMS dust optical depth and of the 10-m wind speed over the
metropolitan cells and by wind sector, and Spearman correlations. The Mann--Whitney test of the
two wind sectors gives $p=0.22$.}
\label{tab:dustwind}
\resizebox{\textwidth}{!}{\begin{tabular}{lrrr}
\toprule
Stratum & Days & Median of stratifier & Metropolitan minus Al~Ain (without Al~Mafraq) \\
\midrule
Dust AOD, low tercile & 108 & 0.03 & $-$58 ($-$52) \\
Dust AOD, mid tercile & 108 & 0.09 & $-$56 ($-$49) \\
Dust AOD, high tercile & 108 & 0.21 & $-$61 ($-$39) \\
\midrule
10-m wind speed (m\,s$^{-1}$), low tercile & 108 & 2.6 & $-$57 ($-$48) \\
10-m wind speed (m\,s$^{-1}$), mid tercile & 108 & 3.3 & $-$66 ($-$52) \\
10-m wind speed (m\,s$^{-1}$), high tercile & 108 & 4.8 & $-$55 ($-$40) \\
\midrule
Wind from 270--360$^{\circ}$ & 224 & -- & $-$58 ($-$48) \\
Wind from other directions & 100 & -- & $-$60 ($-$48) \\
\midrule
Spearman correlation & Days & \multicolumn{2}{r}{$\rho$ (without Al~Mafraq)} \\
\midrule
Difference vs dust AOD & 324 & \multicolumn{2}{r}{$-$0.01 ($+$0.21)} \\
Difference vs wind speed & 324 & \multicolumn{2}{r}{$+$0.06 ($+$0.13)} \\
Regional level vs dust AOD & 324 & \multicolumn{2}{r}{$+$0.64} \\
CAMS dust AOD vs AERONET coarse-mode AOD & 226 & \multicolumn{2}{r}{$+$0.85} \\
CAMS dust share vs AERONET fine-mode fraction & 226 & \multicolumn{2}{r}{$-$0.76} \\
\bottomrule
\end{tabular}
}
\end{table}

\end{document}
